% ============================================================
%  El problema del obstáculo:
%  indicatrices, conos normales, regularización de Moreau–Yosida
%  y método de Newton semisuave
% ------------------------------------------------------------
%  Compilar con pdflatex (dos pasadas) o en Overleaf.
%  Paquetes requeridos: babel-spanish, amsmath, amssymb, amsthm,
%  mathtools, booktabs, geometry, algorithm, algpseudocode, hyperref.
% ============================================================
\documentclass[11pt,a4paper]{article}

\usepackage[utf8]{inputenc}
\usepackage[T1]{fontenc}
\usepackage[spanish,es-noshorthands]{babel}
\usepackage{amsmath,amssymb,amsthm}
\usepackage{mathtools}
\usepackage{booktabs}
\usepackage[margin=2.6cm]{geometry}
\usepackage{algorithm}
\usepackage{algpseudocode}
%\usepackage[colorlinks=true,linkcolor=blue!55!black,citecolor=blue!55!black]{hyperref}
\usepackage{hyperref}

% ----------------- Entornos -----------------
\newtheorem{teorema}{Teorema}[section]
\newtheorem{proposicion}[teorema]{Proposición}
\newtheorem{lema}[teorema]{Lema}
\newtheorem{corolario}[teorema]{Corolario}
\theoremstyle{definition}
\newtheorem{definicion}[teorema]{Definición}
\newtheorem{ejemplo}[teorema]{Ejemplo}
\theoremstyle{remark}
\newtheorem{observacion}[teorema]{Observación}

% ----------------- Notación -----------------
\newcommand{\R}{\mathbb{R}}
\newcommand{\Norm}[1]{\left\lVert #1 \right\rVert}
\DeclareMathOperator{\dist}{dist}
\DeclareMathOperator{\inte}{int}
\newcommand{\dk}{\delta_K}
\newcommand{\dke}{\delta_K^\varepsilon}
\newcommand{\pk}{P_K}

\title{\textbf{El problema del obstáculo:}\\[2mm]
\large indicatrices, conos normales, regularización de Moreau--Yosida\\
y método de Newton semisuave}
\author{Notas de curso}
\date{\today}

\begin{document}
\maketitle

\begin{abstract}
Estas notas desarrollan, con todas las demostraciones, el análisis del problema
del obstáculo: la subdiferencial de la indicatriz de un convexo como cono normal,
la regularización de Moreau--Yosida de esa indicatriz (envolvente, proyección,
penalización explícita y aproximación de Yosida), la convergencia de los problemas
penalizados hacia la solución del problema con restricciones, y finalmente el
método de Newton semisuave para el sistema penalizado, con su interpretación como
método de conjunto activo.
\end{abstract}

\tableofcontents

% ============================================================
\section{Planteamiento del problema}
% ============================================================

Sea $\Omega \subset \R^d$ un dominio acotado y lipschitziano,
$f \in L^2(\Omega)$ y un obstáculo $\psi$ con, por ejemplo,
$\psi \in H^1(\Omega)$ y $\psi \le 0$ sobre $\partial\Omega$
(de modo que $K \neq \emptyset$).
El \emph{problema del obstáculo} consiste en encontrar la deformación
$u^\ast$ de una membrana elástica que está por encima del obstáculo:
\begin{equation}\label{eq:obstaculo}
\boxed{\; u^\ast \in K \text{ y } \int_\Omega \nabla u^\ast \cdot \nabla(v - u^\ast)\,dx
\;\ge\; \int_\Omega f\,(v - u^\ast)\,dx
\quad \forall\, v \in K, \;}
\end{equation}
donde el conjunto de funciones admisibles es el convexo cerrado
\[
K \;=\; \bigl\{ v \in H^1_0(\Omega) \;:\; v \ge \psi \ \text{c.s. en } \Omega \bigr\}.
\]

Con la forma bilineal coerciva $a(u,v) = \int_\Omega \nabla u\cdot\nabla v\,dx$,
el problema \eqref{eq:obstaculo} es exactamente la condición de optimalidad
de primer orden del problema de minimización de energía
\begin{equation}\label{eq:minimizacion}
\min_{u \in H^1_0(\Omega)}\; J(u)
\;=\; \frac12 \int_\Omega |\nabla u|^2\,dx \;-\; \int_\Omega f\,u\,dx
\;+\; \dk(u),
\end{equation}
donde $\dk$ es la \emph{indicatriz} del convexo $K$:
\begin{equation}\label{eq:indicatriz}
\dk(u) \;=\;
\begin{cases}
0 & \text{si } u \in K,\\[1mm]
+\infty & \text{si } u \notin K.
\end{cases}
\end{equation}
La única parte no regular de $J$ es $\dk$: una «función» que vale $0$ sobre
$K$ y salta a $+\infty$ fuera. Toda la dificultad del problema reside en esta
discontinuidad vertical, y toda esta nota consiste en estudiarla y en

\begin{observacion}[Marco funcional]
La restricción $u \ge \psi$ tiene sentido en $L^2(\Omega)$, y
$H^1_0(\Omega) \hookrightarrow L^2(\Omega)$ con inclusión compacta.
Trabajaremos por comodidad en $H = L^2(\Omega)$ cuando estudiemos la
estructura de $K$ y la proyección, y en $V = H^1_0(\Omega)$ para la parte
variacional. El producto escalar de $L^2$ se escribe $(\cdot,\cdot)$ y el de
$H^1_0$ con la forma $a(\cdot,\cdot)$.
\end{observacion}

% ============================================================
\section{La subdiferencial de la indicatriz es el cono normal}
% ============================================================

\begin{definicion}[Subdiferencial convexa]
Sea $F : H \to \R \cup \{+\infty\}$ convexa, propia y cerrada.
La subdiferencial de $F$ en $u$ es
\[
\partial F(u) \;=\; \bigl\{ w \in H \;:\;
F(v) \;\ge\; F(u) + (w,\, v-u) \quad \forall\, v \in H \bigr\}.
\]
Es decir: $w \in \partial F(u)$ si el gráfico de $F$ está por encima del
hiperplano afín $v \mapsto F(u) + (w, v-u)$ tangente en $u$.
\end{definicion}

\begin{definicion}[Cono normal]
Sea $K \subset H$ un convexo cerrado no vacío y $u \in K$. El \emph{cono
normal} a $K$ en $u$ es
\[
N_K(u) \;=\; \bigl\{ w \in H \;:\; (w,\, v - u) \;\le\; 0 \quad \forall\, v \in K \bigr\}.
\]
Geométricamente: $w$ es normal a $K$ en $u$ si forma un ángulo obtuso con
todas las direcciones $v - u$ que apuntan hacia el interior del convexo.
Es un cono cerrado y convexo, y por definición $N_K(u) = \emptyset$ si $u \notin K$.
\end{definicion}

\begin{teorema}[$\partial \dk = N_K$]\label{thm:ind-normal}
Sea $K \subset H$ convexo cerrado no vacío. Entonces, para todo $u \in H$,
\[
\partial \dk(u) \;=\; N_K(u).
\]
\end{teorema}

\begin{proof}
Distinguimos tres casos.

\smallskip
\emph{Caso 1: $u \notin K$.} Entonces $\dk(u) = +\infty$ y, por la convención
usual, $\partial \dk(u) = \emptyset = N_K(u)$.

\smallskip
\emph{Caso 2: $u \in K$.} Probamos la doble inclusión.

($\subseteq$) Sea $w \in \partial\dk(u)$. La desigualdad de subtangente dice
\[
\dk(v) \;\ge\; \dk(u) + (w, v-u) \quad \forall v \in H.
\]
Como $\dk(u) = 0$, si tomamos $v \in K$ (de modo que $\dk(v) = 0$) obtenemos
\[
0 \;\ge\; (w, v - u) \quad \forall\, v \in K,
\]
que es exactamente $w \in N_K(u)$. Nótese que para $v \notin K$ la
desigualdad es trivial ($+\infty \ge$ cualquier cosa), así que no aporta
nada más.

($\supseteq$) Sea $w \in N_K(u)$, es decir $(w, v-u) \le 0$ para todo
$v \in K$. Verifiquemos la desigualdad de subtangiente:
\begin{itemize}
\item si $v \in K$: $\dk(v) = 0 \ge (w, v-u)$ por hipótesis, y
$\dk(u) + (w,v-u) = (w, v-u) \le 0$;\,✓
\item si $v \notin K$: $\dk(v) = +\infty$, y la desigualdad es automática.\,✓
\end{itemize}
Por lo tanto $w \in \partial\dk(u)$.

\smallskip
\emph{Conclusión:} en ambos casos los dos conjuntos coinciden.
\end{proof}

\begin{proposicion}[Puntos interiores]\label{prop:interior}
Si $u \in \inte K$, entonces $N_K(u) = \{0\}$.
\end{proposicion}

\begin{proof}
Sea $w \in N_K(u)$. Como $u$ es interior, existe $\rho > 0$ tal que
$u \pm t w \in K$ para todo $t \in (0, \rho/\Norm{w}]$. Aplicando la
desigualdad del cono normal con $v = u + t w$ y con $v = u - t w$:
\[
(w,\, t w) \le 0 \quad\text{y}\quad (w,\, -t w) \le 0
\qquad\Longrightarrow\qquad
t\,\Norm{w}^2 \le 0 .
\]
Luego $w = 0$.
\end{proof}

\begin{observacion}
La proposición expresa un hecho físico evidente: en el interior del convexo
la restricción no actúa, la reacción de contacto es nula. La reacción
(reacción = elemento del cono normal = multiplicador de Lagrange) solo puede
ser no nula sobre la \emph{frontera} de $K$, es decir, sobre la zona de
contacto $\{u^\ast = \psi\}$. Este es el origen de las condiciones de
complementariedad del problema del obstáculo.
\end{observacion}

\subsection{Monotonía y monotonía maximal de $N_K$}

Recordemos que un operador multivaluado $A : H \rightrightarrows H$ es
\emph{monótono} si
\[
(w - z,\, x - y) \;\ge\; 0
\qquad \forall\, x, y \in H,\ \forall\, w \in A(x),\ z \in A(y),
\]
y es \emph{maximal monótono} si ningún par $(x,w)$ puede añadirse a su
gráfico sin destruir la monotonía.

\begin{proposicion}\label{prop:monotono}
$N_K$ es un operador monótono.
\end{proposicion}

\begin{proof}
Sean $x, y \in K$, $w_x \in N_K(x)$, $w_y \in N_K(y)$. Por definición,
aplicando cada cono normal al punto del otro:
\[
(w_x,\, y - x) \le 0
\qquad\text{(tomando } v = y \text{ en } N_K(x)\text{)},
\]
\[
(w_y,\, x - y) \le 0
\qquad\text{(tomando } v = x \text{ en } N_K(y)\text{)}.
\]
Sumando ambas desigualdades y reordenando:
\[
(w_x - w_y,\, x - y) \;=\; -(w_x, y - x) - (w_y, x - y) \;\ge\; 0. \qedhere
\]
\end{proof}

\begin{teorema}[Maximalidad]\label{thm:maximal}
$N_K$ es maximal monótono en $H$.
\end{teorema}

\begin{proof}
Supongamos que $(x, w)$ se puede añadir al gráfico de $N_K$ conservando la
monotonía, es decir
\begin{equation}\label{eq:hip-max}
(w - z,\, x - u) \;\ge\; 0
\qquad \forall\, u \in K,\ \forall\, z \in N_K(u).
\end{equation}
Debemos probar que $w \in N_K(x)$.

\emph{Paso 1: construimos un buen punto de prueba.}
Para $t > 0$ sea $u_t := \pk(x + t w) \in K$ la proyección de $x + t w$
sobre $K$. La caracterización variacional de la proyección dice
\[
(y - \pk(y),\, v - \pk(y)) \le 0 \quad \forall v \in K,
\]
que aplicada con $y = x + t w$ da
\[
x + t w - u_t \;\in\; N_K(u_t).
\]
Como $N_K(u_t)$ es un cono, también
\[
z_t := \frac{x + t w - u_t}{t} \;\in\; N_K(u_t).
\]

\emph{Paso 2: explotamos la hipótesis.}
Escribiendo \eqref{eq:hip-max} con $u = u_t$ y $z = z_t$:
\[
(w - z_t,\, x - u_t) \;\ge\; 0.
\]
Ahora bien, descomponiendo $x + t w - u_t = (x - u_t) + t w$,
\[
(z_t,\, x - u_t)
= \frac{1}{t}\Norm{x - u_t}^2 + (w,\, x - u_t).
\]
Sustituyendo:
\[
(w,\, x - u_t) - \frac{1}{t}\Norm{x - u_t}^2 - (w,\, x - u_t) \;\ge\; 0
\qquad\Longrightarrow\qquad
\Norm{x - u_t}^2 \;\le\; 0.
\]
¡La desigualdad degeneró! Concluimos que $u_t = x$ para todo $t > 0$.

\emph{Paso 3: conclusión.}
Entonces $x \in K$ (es proyección de un punto, luego pertenece a $K$) y
\[
x + t w - u_t \;=\; t w \;\in\; N_K(x),
\]
y como $N_K(x)$ es un cono, $w \in N_K(x)$. $\blacksquare$
\end{proof}

\begin{observacion}
Este argumento es la demostración general (Debrunner--Haine / Rockafellar)
de que todo $\partial F$ convexo propio cerrado es maximal monótono,
especializada al caso de la indicatriz, donde la resolución del paso
variacional es la proyección. La estructura es siempre la misma: la
hipótesis de monotonía «aplana» la desigualdad y fuerza la conclusión.
\end{observacion}

% ============================================================
\section{Paso 1: la regularización de Moreau--Yosida de $\delta_K$}
% ============================================================

La idea de Moreau es sustituir la pared vertical de la indicatriz por una
parábola de pendiente finita: en lugar de pagar $+\infty$ al salir de $K$,
se paga la distancia al cuadrado.

\begin{definicion}[Envolvente de Moreau]\label{def:envolvente}
Para $\varepsilon > 0$ y $u \in H$,
\[
\dke(u)
\;=\; \min_{v \in H} \left\{ \dk(v) + \frac{1}{2\varepsilon}\Norm{u - v}^2 \right\}.
\]
\end{definicion}

\subsection{Existencia y unicidad del minimizador}

\begin{proposicion}\label{prop:existencia}
Para todo $u \in H$ y todo $\varepsilon > 0$, el problema de la
Definición~\ref{def:envolvente} tiene un único minimizador
$p_\varepsilon(u) \in K$. Además $p_\varepsilon(u) = \pk(u)$, la proyección
de $u$ sobre $K$.
\end{proposicion}

\begin{proof}
\emph{Existencia (método directo).} Sea
$\Phi_u(v) = \dk(v) + \frac{1}{2\varepsilon}\Norm{u-v}^2$.
Desarrollando el cuadrado,
\[
\Phi_u(v)
\;\ge\; \frac{1}{2\varepsilon}\Norm{v}^2
- \frac{1}{\varepsilon}(u, v)
+ \frac{1}{2\varepsilon}\Norm{u}^2 .
\]
Es decir, fuera de $K$ el funcional vale $+\infty$, y sobre $K$ es
coercivo: la parte cuadrática domina la parte lineal por Young,
\[
\frac{1}{2\varepsilon}\Norm{v}^2 - \frac{1}{\varepsilon}(u,v)
\;\ge\; \frac{1}{4\varepsilon}\Norm{v}^2 - \frac{1}{\varepsilon}\Norm{u}^2 ,
\]
así que toda sucesión minimizante está acotada. Tomemos $(v_n) \subset K$
minimizante. Como $K$ es convexo y cerrado, es débilmente secuencialmente
cerrado; por compacidad débil, $v_n \rightharpoonup p \in K$. El funcional
$\Phi_u$ es débilmente semicontinuo inferior (norma y producto escalar lo
son, y $\dk$ es débilmente s.c.i. porque $K$ es débilmente cerrado). Luego
$p$ es minimizador.

\smallskip
\emph{Unicidad.} $\Phi_u$ es estrictamente convexo: la indicatriz de un
convexo es convexa, y $\frac{1}{2\varepsilon}\Norm{u - \cdot}^2$ es
estrictamente convexa. Un funcional estrictamente convexo tiene a lo sumo un
minimizador.

\smallskip
\emph{Identificación con la proyección.} Por la regla de suma de
subdiferenciales (válida sin condición de calificación aquí, porque
$v \mapsto \frac{1}{2\varepsilon}\Norm{u-v}^2$ es de clase $C^1$ en todo
$H$ con gradiente continuo):
\[
0 \in \partial \Phi_u(p)
= \partial \dk(p) + \frac{p - u}{\varepsilon}.
\]
Por el Teorema~\ref{thm:ind-normal}, $\frac{u - p}{\varepsilon} \in N_K(p)$,
es decir
\begin{equation}\label{eq:carac-proy}
(u - p,\, v - p) \;\le\; 0 \qquad \forall\, v \in K,
\end{equation}
que es exactamente la caracterización variacional de la proyección:
$p = \pk(u)$. (Para el lector que no recuerde la equivalencia:
\eqref{eq:carac-proy} es la derivación en $t=0$ del hecho de que
$t \mapsto \Norm{u - \pk(u) + t(p - \pk(u))}^2$ tiene un mínimo en $t=0$
sobre la cuerda de $K$; la implicación inversa se obtiene integrando a lo
largo de la cuerda $p + t(v-p)$, $t\in[0,1]$.)
\end{proof}

\subsection{Valor de la envolvente y primera fórmula fundamental}

\begin{corolario}[Fórmula de la envolvente]\label{cor:valor}
\[
\dke(u) \;=\; \frac{1}{2\varepsilon}\, \Norm{u - \pk(u)}^2
\;=\; \frac{1}{2\varepsilon}\, \dist(u, K)^2 .
\]
\end{corolario}

\begin{proof}
El minimizador $p = \pk(u)$ pertenece a $K$, luego $\dk(p) = 0$ y solo
queda el término cuadrático. La segunda igualdad es la definición de la
distancia a un convexo cerrado, alcanzada en la proyección.
\end{proof}

\begin{observación}[Lectura de la fórmula]
La envolvente de Moreau de la indicatriz es la \emph{penalización
cuadrática} de la distancia al conjunto: $\dke \le \dk$ para todo $u$,
con igualdad exactamente en $K$ (donde ambas valen $0$). La pared vertical
se ha redondeado en un valle parabólico de anchura $O(\sqrt{\varepsilon})$
alrededor del convexo.
\end{observación}

\begin{corolario}[Resolvente de $N_K$]\label{cor:resolvente}
Para todo $\lambda > 0$,
\[
J_\lambda^{N_K}(u) := (I + \lambda N_K)^{-1}(u) = \pk(u),
\]
y $J_\lambda^{N_K}$ es unívoco y contractante.
\end{corolario}

\begin{proof}
$w \in (I + \lambda N_K)^{-1}(u)$ significa $u \in w + \lambda N_K(w)$, es
decir $\frac{u - w}{\lambda} \in N_K(w)$, que es
\eqref{eq:carac-proy} con $p = w$. Luego $w = \pk(u)$, único por la
Proposición~\ref{prop:existencia}. La contractividad es la no-expansividad
de la proyección:
\[
\Norm{\pk(u) - \pk(u')} \;\le\; \Norm{u - u'} ,
\]
que se obtiene sumando $(u - \pk(u), v - \pk(u)) \le 0$ con
$v = \pk(u')$ y la desigualdad simétrica, y aplicando Cauchy--Schwarz.
\end{proof}

\begin{observación}
Este corolario es el teorema de Minty hecho carne: el operador maximal
monótono multivaluado $N_K$ tiene resolvente unívoco y lipschitziano, y ese
resolvente es la proyección. Es el punto exacto donde la teoría abstracta
de la sección anterior se convierte en un objeto computable.
\end{observación}

\subsection{Diferenciabilidad de la envolvente y aproximación de Yosida}

\begin{teorema}\label{thm:diferenciable}
La envolvente $\dke$ es Fréchet-diferenciable en todo $H$, y
\[
\nabla \dke(u) \;=\; \frac{u - \pk(u)}{\varepsilon}.
\]
\end{teorema}

\begin{proof}
Fijemos $u \in H$, escribamos $p = \pk(u)$ y, para $h \in H$,
$p_h = \pk(u + h)$. Recordemos que la proyección es no expansiva:
$\Norm{p - p_h} \le \Norm{h}$.

\smallskip
\emph{Cota superior.} Probando $v = p$ en el problema que define
$\dke(u+h)$:
\[
\dke(u+h) \;\le\; \frac{1}{2\varepsilon}\Norm{u + h - p}^2
\;=\; \dke(u) + \frac{1}{\varepsilon}(u - p,\, h) + \frac{1}{2\varepsilon}\Norm{h}^2 ,
\]
donde el último paso es la identidad polar
$\Norm{(u-p)+h}^2 = \Norm{u-p}^2 + 2(u-p,h) + \Norm{h}^2$ y el
Corolario~\ref{cor:valor}.

\smallskip
\emph{Cota inferior.} Como $p_h$ minimiza en $u + h$ y $p_h \in K$, el
Corolario~\ref{cor:valor} da
$\dke(u+h) = \frac{1}{2\varepsilon}\Norm{u + h - p_h}^2$, mientras que
probando $v = p_h$ en el problema que define $\dke(u)$:
\[
\dke(u) \;\le\; \frac{1}{2\varepsilon}\Norm{u - p_h}^2 .
\]
Restando:
\[
\dke(u+h) - \dke(u)
\;\ge\; \frac{1}{2\varepsilon}\Bigl(\Norm{u + h - p_h}^2 - \Norm{u - p_h}^2\Bigr)
\;=\; \frac{1}{\varepsilon}(u - p_h,\, h) + \frac{1}{2\varepsilon}\Norm{h}^2 .
\]

\smallskip
\emph{Conclusión.} Combinando ambas cotas con el gradiente candidato
$g = \frac{u - p}{\varepsilon}$:
\begin{align*}
\text{(inferior)}\quad
\dke(u+h) - \dke(u) - (g, h)
&\;\ge\; \frac{1}{\varepsilon}(p - p_h,\, h) + \frac{1}{2\varepsilon}\Norm{h}^2
\;\ge\; -\frac{1}{\varepsilon}\Norm{h}^2 + \frac{1}{2\varepsilon}\Norm{h}^2
\;=\; -\frac{1}{2\varepsilon}\Norm{h}^2,\\[1mm]
\text{(superior)}\quad
\dke(u+h) - \dke(u) - (g, h)
&\;\le\; \frac{1}{2\varepsilon}\Norm{h}^2 .
\end{align*}
donde hemos usado Cauchy--Schwarz y la no-expansividad en
$(p - p_h, h) \ge -\Norm{p-p_h}\,\Norm{h} \ge -\Norm{h}^2$. Por lo tanto
\[
\Bigl|\, \dke(u+h) - \dke(u) - \tfrac{1}{\varepsilon}(u - \pk(u),\, h) \,\Bigr|
\;\le\; \frac{1}{2\varepsilon}\Norm{h}^2 \;\xrightarrow[h \to 0]{}\; 0 ,
\]
lo que es exactamente la diferenciabilidad Fréchet con el gradiente
anunciado (¡y con resto explícito $O(\Norm{h}^2)$!).
\end{proof}

\begin{definicion}[Aproximación de Yosida]
El operador
\[
(N_K)_\varepsilon
\;:=\; \frac{1}{\varepsilon}\bigl(I - \pk\bigr)
\;=\; \nabla \dke
\]
se llama \emph{aproximación de Yosida} de $N_K$.
\end{definicion}

\begin{proposicion}\label{prop:yosida-prop}
$(N_K)_\varepsilon$ es:
\begin{enumerate}
\item unívoco y $\tfrac{1}{\varepsilon}$-lipschitziano;
\item monótono (es más: cocoercivo, $\varepsilon\,\Norm{(N_K)_\varepsilon u - (N_K)_\varepsilon v}^2 \le \bigl((N_K)_\varepsilon u - (N_K)_\varepsilon v,\, u - v\bigr)$);
\item maximal monótono;
\item exacto sobre $K^\ast := \{ u : 0 \in N_K(u) \}$:
$(N_K)_\varepsilon(u) = 0 \iff u \in K$ con $N_K(u) \ni 0$, es decir
$(N_K)_\varepsilon(u) = 0 \iff u = \pk(u) \iff u \in K$.
\end{enumerate}
\end{proposicion}

\begin{proof}
(1) es la no-expansividad de $\pk$. Para (2), usando la caracterización
variacional de la proyección dos veces
$(u - \pk(u),\, \pk(v) - \pk(u)) \le 0$ y la simétrica:
\begin{align*}
\bigl((N_K)_\varepsilon u - (N_K)_\varepsilon v,\; u - v\bigr)
&= \frac{1}{\varepsilon}\Bigl(\Norm{u-v}^2 - \bigl(\pk(u) - \pk(v),\, u-v\bigr)\Bigr)\\
&\ge\; \frac{1}{\varepsilon}\Bigl(\Norm{u-v}^2 - \Norm{\pk(u)-\pk(v)}\,\Norm{u-v}\Bigr)
\;\ge\; 0 ,
\end{align*}
de hecho $\ge \frac{1}{\varepsilon}\Norm{\pk(u)-\pk(v)}^2 \ge 0$ si se usa
la firme no-expansividad $\Norm{p_u - p_v}^2 \le (u-v, p_u - p_v)$.
(3) Un gradiente de funcional convexo diferenciable definido en todo $H$
es maximal monótono (teorema de Rockafellar). (4) es inmediato de
$(N_K)_\varepsilon(u) = \frac{1}{\varepsilon}(u - \pk(u))$.
\end{proof}

% ============================================================
\section{Paso 2: el cálculo explícito para el obstáculo}
% ============================================================

Hasta ahora el desarrollo vale para \emph{todo} convexo cerrado. Ahora
explotamos la estructura puntual del conjunto del obstáculo
\[
K = \bigl\{ v \in L^2(\Omega) : v \ge \psi \ \text{c.s.}\bigr\}.
\]

\begin{teorema}[Proyección puntual sobre el conjunto del obstáculo]\label{thm:proy-puntual}
Para todo $u \in L^2(\Omega)$,
\[
\pk(u)(x) \;=\; \max\bigl(u(x),\, \psi(x)\bigr) \quad \text{c.s. en } \Omega,
\]
y por lo tanto
\[
u - \pk(u) \;=\; (u - \psi)_- ,\qquad
\Norm{u - \pk(u)}_{L^2}^2 \;=\; \int_\Omega \bigl(u - \psi\bigr)_-^2\,dx .
\]
\end{teorema}

\begin{proof}
Escribamos $w := \max(u, \psi)$ y verifiquemos que $w$ satisface la
caracterización variacional \eqref{eq:carac-proy}; por unicidad de la
proyección (Proposición~\ref{prop:existencia}) concluimos.

\smallskip
\emph{(i) $w \in K$:} evidente, $\max(u,\psi) \ge \psi$ c.s.

\smallskip
\emph{(ii) $u - w = (u - \psi)_-$:} en cada punto, si $u(x) \ge \psi(x)$
entonces $w(x) = u(x)$ y $u - w = 0 = (u-\psi)_-(x)$; si $u(x) < \psi(x)$
entonces $w(x) = \psi(x)$ y $u - w = u - \psi = (u-\psi)_-(x) < 0$. En
ambos casos $u - w = \min(u - \psi, 0) = (u-\psi)_-$. En particular
\[
\Norm{u - w}_{L^2}^2 = \int_\Omega (u-\psi)_-^2\,dx .
\]

\emph{(iii) Desacoplamiento puntual:} el funcional
$v \mapsto \int_\Omega (u(x) - v(x))^2\,dx$ sujeto a $v \ge \psi$ es una
integral de funcionales puntuales independientes en cada $x$; su
minimizador se obtiene minimizando en cada punto
$v(x) \mapsto (u(x) - v(x))^2$ bajo $v(x) \ge \psi(x)$, cuyo mínimo
evidente es $v(x) = \max(u(x), \psi(x))$. Formalmente basta verificar la
caracterización variacional:

\smallskip
\emph{(iv) $(u - w,\, v - w)_{L^2} \le 0$ para todo $v \in K$:}
descompongamos el dominio en
\[
\Omega_+ = \{ u \ge \psi \}
\qquad\text{y}\qquad
\Omega_- = \{ u < \psi \}.
\]
Sobre $\Omega_+$, $u - w = 0$ y el integrando es nulo. Sobre $\Omega_-$,
$u - w = u - \psi < 0$ y $v - w = v - \psi \ge 0$ (porque $v \in K$).
Luego el integrando
\[
(u - w)(v - w) = \underbrace{(u - \psi)}_{<\,0}\,\underbrace{(v - \psi)}_{\ge\,0}
\;\le\; 0 \quad\text{c.s. en } \Omega_- ,
\]
y por integración $(u-w,\,v-w)_{L^2} \le 0$. $\blacksquare$
\end{proof}

\begin{corolario}[Fórmulas explícitas para el obstáculo]\label{cor:explicito}
Para todo $u \in L^2(\Omega)$:
\[
\dke(u) = \frac{1}{2\varepsilon}\int_\Omega (u - \psi)_-^2\,dx,
\qquad
\nabla\dke(u) = \frac{1}{\varepsilon}(u - \psi)_-,
\qquad
(N_K)_\varepsilon(u) = \frac{1}{\varepsilon}(u - \psi)_- .
\]
\end{corolario}

\begin{proof}
Combinar el Corolario~\ref{cor:valor}, el Teorema~\ref{thm:diferenciable}
y el Teorema~\ref{thm:proy-puntual}.
\end{proof}

\begin{observación}[Interpretación física]
El funcional penalizado
\[
J_\varepsilon(u)
\;=\; \frac12 \int_\Omega |\nabla u|^2\,dx
\;-\; \int_\Omega f\,u\,dx
\;+\; \frac{1}{2\varepsilon}\int_\Omega \bigl(u - \psi\bigr)_-^2\,dx
\]
describe la membrana sujeta por \emph{muelles elásticos} de rigidez
$1/\varepsilon$ distribuidos por todo el dominio, que solo actúan allí
donde la membrana tiende a penetrar el obstáculo. No hay restricción dura:
la reacción de contacto $(1/\varepsilon)(u-\psi)_-$ crece suavemente con la
penetración. Cuando $\varepsilon \to 0$, los muelles se vuelven
infinitamente rígidos y se recupera la restricción.
\end{observación}

% ============================================================
\section{El problema penalizado y su convergencia}
% ============================================================

\subsection{Ecuación de Euler del problema penalizado}

$J_\varepsilon$ es coercivo, estrictamente convexo y continuo en
$H^1_0(\Omega)$ (la penalidad es $C^1$ por el Teorema~\ref{thm:diferenciable},
con gradiente $\frac{1}{\varepsilon}(u - \psi)_- \in L^2 \hookrightarrow
H^{-1}$). Tiene pues un único minimizador $u_\varepsilon$,
caracterizado por la ecuación de Euler no lineal:
\begin{equation}\label{eq:euler-penalizado}
a(u_\varepsilon,\, v) \;+\; \frac{1}{\varepsilon}\bigl( (u_\varepsilon - \psi)_- ,\, v \bigr)_{L^2}
\;=\; \int_\Omega f\,v\,dx
\qquad \forall\, v \in H^1_0(\Omega).
\end{equation}
La violación de la restricción se controla \emph{en una línea}:

\begin{proposicion}[Violación de la restricción]\label{prop:violacion}
Existe $C > 0$ independiente de $\varepsilon$ tal que
\[
\Norm{(u_\varepsilon - \psi)_-}_{L^2} \;\le\; \sqrt{2 C \,\varepsilon} .
\]
\end{proposicion}

\begin{proof}
Como $u^\ast \in K$ es admisible para $J_\varepsilon$ con
$J_\varepsilon(u^\ast) = J(u^\ast)$ (su penalidad es nula), y como
$J_\varepsilon(u) \ge \frac12\Norm{\nabla u}^2 - (f, u) \ge
-\frac12 \Norm{f}_{H^{-1}}^2$ para todo $u$ (completando el cuadrado):
\[
\frac{1}{2\varepsilon}\Norm{(u_\varepsilon - \psi)_-}_{L^2}^2
\;\le\; J_\varepsilon(u_\varepsilon) \;\le\; J_\varepsilon(u^\ast) = J(u^\ast)
\;\le\; J(u^\ast) + \frac12\Norm{f}_{H^{-1}}^2 .
\]
\end{proof}

\subsection{Convergencia $u_\varepsilon \to u^\ast$}

\begin{teorema}\label{thm:convergencia}
Cuando $\varepsilon \downarrow 0$,
\[
u_\varepsilon \rightharpoonup u^\ast \text{ débilmente en } H^1_0,
\qquad
u_\varepsilon \to u^\ast \text{ fuertemente en } L^2,
\]
donde $u^\ast$ es la solución del problema del obstáculo \eqref{eq:obstaculo}.
\end{teorema}

\begin{proof}
\emph{Paso 1: acotación.} De $J_\varepsilon(u_\varepsilon) \le J(u^\ast)$ y
$J_\varepsilon \ge \frac12\Norm{\nabla\cdot}^2 - (f,\cdot)$ deducimos
$\Norm{u_\varepsilon}_{H^1_0} \le C$ uniformemente.

\emph{Paso 2: paso al límite en la restricción.}
Por la Proposición~\ref{prop:violacion}, $(u_\varepsilon - \psi)_- \to 0$
fuertemente en $L^2$. Por compacidad de Rellich, de toda sucesión
$\varepsilon_n \to 0$ se extrae $u_{\varepsilon_n} \rightharpoonup \bar u$
débilmente en $H^1_0$ y fuertemente en $L^2$. Entonces
$(u_{\varepsilon_n} - \psi)_- \to (\bar u - \psi)_-$ fuertemente en $L^2$
(el operador $(\cdot)_-$ es lipschitziano), luego $(\bar u - \psi)_- = 0$,
es decir $\bar u \in K$.

\emph{Paso 3: $\bar u$ resuelve el problema.} Semi-continuidad inferior
débil de la norma y convergencia fuerte del término lineal dan
\[
J(\bar u) \;\le\; \liminf J_\varepsilon(u_\varepsilon)
\;\le\; J(u^\ast),
\]
donde la primera desigualdad usa que la penalidad tiende a $0$
(Proposición~\ref{prop:violacion}) y que $\delta_K(\bar u) = 0$. Como
$u^\ast$ es el único minimizador de $J$ sobre $K$, $\bar u = u^\ast$.
La unicidad del punto límite débil implica la convergencia de toda la
familia.
\end{proof}

\begin{teorema}[Cota en norma de energía; Brezis--Stampacchia]
\label{thm:BS}
Si además $f \in L^2(\Omega)$, $\psi$ es suficientemente regular
(por ejemplo $\psi \in H^2$ con $\psi \le 0$ en $\partial\Omega$) y
$u^\ast \in H^2(\Omega)$, entonces existe $C$ independiente de
$\varepsilon$ tal que
\[
\Norm{u_\varepsilon - u^\ast}_{H^1_0} \;\le\; C\,\sqrt{\varepsilon}.
\]
\end{teorema}

\begin{proof}[Idea de la demostración]
Se resta la ecuación de Euler \eqref{eq:euler-penalizado} de la ecuación del
obstáculo $a(u^\ast, v) + \langle \lambda^\ast, v\rangle = (f, v)$
(que vale para todo $v \in H^1_0$ por densidad, con el multiplicador
$\lambda^\ast \in N_K(u^\ast) \subset H^{-1}$) con $v = u_\varepsilon -
u^\ast$, y se acota el término
$\langle (N_K)_\varepsilon(u_\varepsilon) - \lambda^\ast,\,
u_\varepsilon - u^\ast\rangle$ descomponiendo
$u_\varepsilon - u^\ast = (u_\varepsilon - \psi)_- + (z_\varepsilon - u^\ast)$
con $z_\varepsilon := \max(u_\varepsilon, \psi) \in K$ y usando:
$\langle \lambda^\ast, z_\varepsilon - u^\ast\rangle \le 0$ (cono normal),
$(\lambda^\ast, (u_\varepsilon - \psi)_-) \ge 0$ en el sentido de medidas
($\lambda^\ast \ge 0$ y $(u_\varepsilon-\psi)_- \le 0$), y la
Proposición~\ref{prop:violacion}. Se obtiene
$\Norm{u_\varepsilon - u^\ast}_{H^1_0}^2 \le C'\, \varepsilon$.
Para la demostración completa, ver \cite{BS68, Glowinski}.
\end{proof}

\subsection{Convergencia del multiplicador}

Definimos la reacción de los muelles
\[
\lambda_\varepsilon := \frac{1}{\varepsilon}(u_\varepsilon - \psi)_-
\qquad (\text{un elemento de } H^{-1} \text{ vía }
\langle \lambda_\varepsilon, v\rangle := (\lambda_\varepsilon, v)_{L^2}).
\]

\begin{teorema}\label{thm:multiplicador}
De toda sucesión $\varepsilon_n \to 0$ se puede extraer una subsucesión tal
que $\lambda_{\varepsilon_n} \overset{\ast}{\rightharpoonup} \lambda$
débilmente-$\ast$ en $H^{-1}$, donde $\lambda$ satisface
\[
-\Delta u^\ast + \lambda = f \ \text{en } H^{-1},
\qquad
\lambda \in N_K(u^\ast).
\]
En consecuencia, en el límite se recuperan las condiciones de
complementariedad del problema del obstáculo:
\[
u^\ast \ge \psi, \qquad \lambda \ge 0 \text{ (sobre la zona de contacto)},
\qquad \lambda = 0 \text{ fuera del contacto}.
\]
\end{teorema}

\begin{proof}
\emph{Acotación en $H^{-1}$:} de \eqref{eq:euler-penalizado},
\[
\langle \lambda_\varepsilon, v \rangle = (f, v) - a(u_\varepsilon, v)
\qquad\Longrightarrow\qquad
\Norm{\lambda_\varepsilon}_{H^{-1}} \le \Norm{f}_{H^{-1}} + C\,\Norm{u_\varepsilon}_{H^1_0}.
\]
Extraemos una subsucesión débilmente-$\ast$ convergente a $\lambda$.
Pasando al límite en la ecuación anterior (con
$u_{\varepsilon_n} \rightharpoonup u^\ast$) se obtiene
$-\Delta u^\ast + \lambda = f$.

\emph{$\lambda \in N_K(u^\ast)$:} sea $v \in K$ arbitrario. Escribamos
$z_\varepsilon := \max(u_\varepsilon, \psi) = u_\varepsilon -
(u_\varepsilon - \psi)_- \in K$. Puntualmente,
$(u_\varepsilon - \psi)_- \le 0$ y $v - z_\varepsilon = v - \max(u_\varepsilon, \psi) \ge 0$
(pues $v \ge \psi$), luego
\[
\langle \lambda_\varepsilon, v \rangle
= \frac{1}{\varepsilon}\bigl( (u_\varepsilon - \psi)_- ,\, v \bigr)
\;\le\; \frac{1}{\varepsilon}\bigl( (u_\varepsilon - \psi)_- ,\, z_\varepsilon \bigr)
= \langle \lambda_\varepsilon, u_\varepsilon \rangle
- \Norm{\lambda_\varepsilon}_{L^2}^2\,\varepsilon
\;\le\; \langle \lambda_\varepsilon, u_\varepsilon \rangle .
\]
Ahora, $\langle \lambda_\varepsilon, u_\varepsilon \rangle
= (f, u_\varepsilon) - a(u_\varepsilon, u_\varepsilon)$ converge a
$(f, u^\ast) - a(u^\ast, u^\ast) = \langle \lambda, u^\ast \rangle$
(convergencia débil en $H^1_0$ para el término lineal,
semi-continuidad inferior para el cuadrático y la ecuación límite para
identificar el valor). Pasando al límite:
\[
\langle \lambda, v \rangle \;\le\; \langle \lambda, u^\ast \rangle
\qquad \forall\, v \in K,
\]
que es la definición de $\lambda \in N_K(u^\ast)$.
\end{proof}

% ============================================================
\section{El método de Newton semisuave para el problema penalizado}
% ============================================================

\subsection{Discretización}

Tras una discretización por elementos finitos $P_1$ (o por diferencias
finitas), el problema penalizado \eqref{eq:euler-penalizado} se convierte
en un sistema no lineal en $\R^N$, con $N$ el número de nodos interiores:
\begin{equation}\label{eq:sistema}
F(u) \;=\; A\,u \;-\; b \;+\; \frac{1}{\varepsilon}\,(u - \psi)_- \;=\; 0,
\end{equation}
donde:
\begin{itemize}
\item $A \in \R^{N\times N}$ es la matriz de rigidez (definida positiva,
SPD tras las condiciones de Dirichlet),
\item $b$ es el vector de carga $(b_i = \int_\Omega f\,\varphi_i\,dx)$,
\item $\psi \in \R^N$ contiene los valores nodales del obstáculo,
\item $(u - \psi)_-$ se aplica \emph{componente a componente}:
$((u-\psi)_-)_i = \min(u_i - \psi_i,\, 0)$.
\end{itemize}

\subsection{Semisuavidad del sistema}

\begin{definicion}[Función semisuave en dimensión finita]
Sea $G : \R^N \to \R^N$ localmente lipschitziana y $x \in \R^N$. $G$ es
\emph{semisuave} en $x$ si
\[
\lim_{h \to 0}\;
\sup_{V \in \partial_C G(x + t h),\; 0 < t \le 1}
\frac{\Norm{G(x+h) - G(x) - V h}}{\Norm{h}} \;=\; 0,
\]
donde $\partial_C$ es el jacobiano generalizado de Clarke. Es
\emph{fuertemente semisuave} si el límite es $O(\Norm{h})$, es decir
$\Norm{G(x+h) - G(x) - Vh} = O(\Norm{h}^2)$.
\end{definicion}

\begin{lema}\label{lema:g-semisuave}
La función escalar $g(t) = \min(t, 0) = -\max(-t, 0)$ es fuertemente
semisuave en todo $\R$, con subdiferencial de Clarke
\[
\partial g(t) \;=\;
\begin{cases}
\{1\} & \text{si } t < 0,\\
\{0\} & \text{si } t > 0,\\
[0, 1] & \text{si } t = 0.
\end{cases}
\]
\end{lema}

\begin{proof}
$g$ es localmente lipschitziana (constante $1$) y suave salvo en $t=0$,
con lo que el cálculo de $\partial g$ es elemental. Para la semisuavidez,
el punto clave es que en la definición el jacobiano $V$ se evalúa
\emph{sobre el segmento} $x + th$, $0 < t \le 1$.

Fijemos $t_0 \in \R$ y $h \neq 0$ pequeño.
\begin{itemize}
\item Si $t_0 > 0$: para $\Norm{h}$ pequeño, todo el segmento
$[t_0, t_0 + h]$ está en $(0, \infty)$, donde $g \equiv$ lineal con
derivada... más precisamente $g = 0$ allí: $g(t_0 + h) - g(t_0) = 0$ y
$\partial g \equiv \{0\}$ sobre el segmento, luego $V = 0$ y el error es
\emph{exactamente nulo}.
\item Si $t_0 < 0$: el segmento está en $(-\infty, 0)$, donde $g(t) = t$:
$g(t_0 + h) - g(t_0) = h$ y $V = 1$, luego $Vh = h$: error nulo.
\item Si $t_0 = 0$: si $h > 0$, el segmento $(0, h]$ está en
$(0,\infty)$, luego $V = 0$ y $g(h) - g(0) - Vh = \min(h,0) - 0 = 0$.
Si $h < 0$, el segmento está en $(-\infty, 0)$, luego $V = 1$ y
$g(h) - g(0) - Vh = h - h = 0$.
\end{itemize}
En todos los casos el error de linealización es idénticamente nulo: $g$
es fuertemente semisuave (de hecho es lineal a trozos).
\end{proof}

\begin{teorema}\label{thm:F-semisuave}
El funcional $F$ de \eqref{eq:sistema} es fuertemente semisuave en
$\R^N$, y su jacobiano generalizado de Clarke es
\[
\partial_C F(u) \;=\; \left\{
A + \frac{1}{\varepsilon}\,D
\;:\;
D = \operatorname{diag}(d_1, \dots, d_N),\;
d_i \in
\begin{cases}
\{1\} & \text{si } u_i < \psi_i,\\
\{0\} & \text{si } u_i > \psi_i,\\
[0,1] & \text{si } u_i = \psi_i
\end{cases}
\right\}.
\]
Todo elemento $V \in \partial_C F(u)$ es definido positivo, y en
particular invertible.
\end{teorema}

\begin{proof}
$F$ es la suma de la aplicación afín $u \mapsto Au - b$ (suave) y de
$u \mapsto \frac{1}{\varepsilon} g(u - \psi)$ con $g$ aplicada
componente a componente. Por el Lema~\ref{lema:g-semisuave} y la regla de
la cadena para funciones semisuaves (la aplicación diagonal
componente a componente preserva la semisuavidad fuerte,
\cite{FacchineiPang}), $F$ es fuertemente semisuave.

Para el jacobiano de Clarke: en un punto $u$, cada componente
$F_i$ solo depende de $u_i$ (más el término afín), luego $\partial_C F(u)$
es el conjunto de las matrices $A + \frac{1}{\varepsilon}D$ con $D$
diagonal y $d_i \in \partial g(u_i - \psi_i)$, exactamente como se
anuncia.

Definición positiva: $A \succ 0$ y $\frac{1}{\varepsilon}D \succeq 0$
(diagonal con entradas $\ge 0$), luego
$V = A + \frac{1}{\varepsilon}D \succeq A \succ 0$ y $\det V \ge \det A > 0$.
\end{proof}

\begin{observación}[Qué significa $\partial_C F$ aquí]
Cada elemento de $\partial_C F(u)$ corresponde a una \emph{selección} de
las derivadas de las piezas lineales: el nodo $i$ está en
\emph{contacto activo} ($u_i < \psi_i$) y contribuye con rigidez
$1/\varepsilon$, está \emph{libre} ($u_i > \psi_i$) y no contribuye, o
está \emph{exactamente en el obstáculo} ($u_i = \psi_i$) y se le puede
asignar cualquier rigidez intermedia. La linealización de Newton es pues
la ecuación lineal de la membrana con un conjunto de resortes «congelado»
en el estado actual.
\end{observacion}

\subsection{El algoritmo}

El método de Newton semisuave aplicado a \eqref{eq:sistema} es:

\begin{algorithm}[H]
\caption{Newton semisuave para el problema del obstáculo penalizado}
\label{alg:newton}
\begin{algorithmic}[1]
\Require $\varepsilon > 0$, vector inicial $u^{(0)} \in \R^N$,
tolerancia $\mathsf{tol} > 0$
\For{$k = 0, 1, 2, \dots$}
  \State Calcular el residuo
  $r^{(k)} := F(u^{(k)}) = A u^{(k)} - b + \frac{1}{\varepsilon}(u^{(k)} - \psi)_-$
  \If{$\Norm{r^{(k)}} \le \mathsf{tol}$}
    \State \textbf{parar}: devolver $u^{(k)}$
  \EndIf
  \State Formar el \emph{conjunto activo} y el \emph{conjunto inactivo}
  \[
  \mathcal{A}_k := \{ i : u_i^{(k)} < \psi_i \},
  \qquad
  \mathcal{I}_k := \{ i : u_i^{(k)} \ge \psi_i \}
  \]
  \State Formar la matriz de Newton
  $V_k := A + \frac{1}{\varepsilon}\,\operatorname{diag}(d^{(k)})$
  con $d_i^{(k)} = 1$ si $i \in \mathcal{A}_k$, $d_i^{(k)} = 0$ si no
  \State Resolver $V_k\, \delta^{(k)} = -\,r^{(k)}$
  \Comment{SPD: factorización de Cholesky}
  \State (Opcional) Busca de línea: elegir $\alpha_k \approx 1$ tal que
  $\Norm{F(u^{(k)} + \alpha_k \delta^{(k)})} < \Norm{r^{(k)}}$
  \State $u^{(k+1)} := u^{(k)} + \alpha_k\,\delta^{(k)}$
\EndFor
\end{algorithmic}
\end{algorithm}

\begin{observación}[Interpretación como conjunto activo]
En cada iteración el método clasifica los nodos en contactantes
($\mathcal{A}_k$) y libres ($\mathcal{I}_k$). Sobre los nodos libres el
sistema de Newton es la ecuación de Poisson; sobre los nodos activos añade
un término de rigidez $1/\varepsilon$ que empuja $u_i$ hacia el obstáculo
(cuanto más pequeño es $\varepsilon$, más se parece a imponer $u_i =
\psi_i$ por penalización). Con $\varepsilon \to 0$ el método tiende al
método \emph{primal--dual de conjunto activo} de Hinterm\"uller--Ito--Kunisch
\cite{HIK02}, que resulta ser precisamente un Newton semisuave sobre la
formulación de complementariedad del problema con restricciones. En
práctica, con complementariedad estricta, el conjunto activo se estabiliza
en un número finito de iteraciones y el método termina.
\end{observacion}

\subsection{Convergencia local cuadrática}

\begin{teorema}\label{thm:newton-conv}
Sea $u^\ast_\varepsilon$ la solución de \eqref{eq:sistema}. Existen
$\rho > 0$ y $C > 0$ tales que si $\Norm{u^{(0)} - u^\ast_\varepsilon}
\le \rho$, la sucesión generada por el Algoritmo~\ref{alg:newton}
(con $\alpha_k = 1$) está bien definida y converge hacia
$u^\ast_\varepsilon$ con velocidad cuadrática:
\[
\Norm{u^{(k+1)} - u^\ast_\varepsilon}
\;\le\; C\,\Norm{u^{(k)} - u^\ast_\varepsilon}^2 .
\]
\end{teorema}

\begin{proof}
Es el teorema estándar de Newton semisuave, cuyas hipótesis verificamos:
\begin{enumerate}
\item \emph{$F$ fuertemente semisuave} en un entorno de
$u^\ast_\varepsilon$: Teorema~\ref{thm:F-semisuave}. Esto da, para
$\Norm{h}$ pequeño y todo $V \in \partial_C F(u^\ast_\varepsilon + h)$:
\[
F(u^\ast_\varepsilon + h) - \underbrace{F(u^\ast_\varepsilon)}_{=0}
- V h \;=\; O(\Norm{h}^2).
\]
\item \emph{Invertibilidad uniforme de los jacobianos:} todo
$V \in \partial_C F(u)$ es SPD
(Teorema~\ref{thm:F-semisuave}); sobre una bola compacta los valores
propios están uniformemente acotados inferiormente por $\lambda_{\min}(A) > 0$,
luego $\Norm{V^{-1}} \le 1/\lambda_{\min}(A)$ uniformemente.
\end{enumerate}
Con estas dos propiedades, el argumento clásico se cierra: si
$e_k := u^{(k)} - u^\ast_\varepsilon$ y $V_k \in \partial_C
F(u^{(k)})$,
\[
e_{k+1}
= V_k^{-1}\Bigl(V_k e_k - F(u^{(k)}) + F(u^\ast_\varepsilon)\Bigr)
= V_k^{-1}\Bigl(F(u^\ast_\varepsilon) - F(u^{(k)}) + V_k\,(u^{(k)} - u^\ast_\varepsilon)\Bigr)
= O(\Norm{e_k}^2),
\]
donde la última igualdad es exactamente la propiedad de semisuavidad fuerte
aplicada con $h = -e_k$ y $V$ evaluado sobre el segmento entre
$u^{(k)}$ y $u^\ast_\varepsilon$ (que es la selección que usa el
algoritmo, con posible redefinición de $d_i \in [0,1]$ en nodos con
$u_i = \psi_i$; todos esos elementos están en $\partial_C F$).
\end{proof}

\begin{observacion}[Sobre la elección de $\varepsilon$ y el condicionamiento]
\leavevmode
\begin{itemize}
\item El error de penalización es $O(\sqrt{\varepsilon})$ en norma $H^1$
(Teorema~\ref{thm:BS}) y $O(\sqrt{\varepsilon})$ en la violación de la
restricción (Proposición~\ref{prop:violacion}): se quiere $\varepsilon$
\emph{pequeño}.
\item El número de condición de las matrices de Newton es
$\kappa(V_k) = O(1/\varepsilon)$: se quiere $\varepsilon$ \emph{grande}.
\item Práctica recomendada: \emph{continuación} $\varepsilon_0 >
\varepsilon_1 > \dots \to \varepsilon_{\min}$, usando la solución anterior
como inicialización. La convergencia cuadrática local hace que cada etapa
cueste pocas iteraciones de Newton.
\item La búsqueda de línea (o un criterio de disminución de
$\Norm{F}$) globaliza el método; sin ella, la convergencia solo está
garantizada localmente.
\end{itemize}
\end{observacion}

\begin{observacion}[Newton semisuave en espacios de Hilbert]
Todo el desarrollo se puede hacer \emph{sin discretizar}: el operador
$B_\varepsilon(u) = \frac{1}{\varepsilon}(u - \psi)_-$ es semisuave como
aplicación $L^p(\Omega) \to L^q(\Omega)$ con $p > q$
(Hinterm\"uller--Ito--Kunisch \cite{HIK02}, Ulbrich \cite{Ulbrich}), y el
Newton semisuave se define directamente en dimensión infinita, sustituyendo
el jacobiano generalizado de Clarke por derivadas newtonianas. La
estructura del algoritmo es idéntica; la discretización
solo interviene al resolver el sistema lineal.
\end{observacion}

% ============================================================
\section{Síntesis}
% ============================================================

\begin{center}
\renewcommand{\arraystretch}{1.35}
\begin{tabular}{@{}ll@{}}
\toprule
\textbf{Objeto} & \textbf{Fórmula} \\
\midrule
Subdiferencial de la indicatriz &
$\partial \delta_K(u) = N_K(u) = \{ w : (w, v-u) \le 0 \ \forall v \in K\}$ \\
Maximalidad & $N_K$ maximal monótono; $J_\lambda^{N_K} = P_K$ (Minty) \\
Envolvente de Moreau &
$\delta_K^\varepsilon(u) = \frac{1}{2\varepsilon}\dist(u,K)^2$ \\
Proyección (obstáculo) &
$P_K(u)(x) = \max(u(x), \psi(x))$, \quad $u - P_K(u) = (u-\psi)_-$ \\
Penalización explícita &
$J_\varepsilon = J + \frac{1}{2\varepsilon}\int (u-\psi)_-^2$ \\
Gradiente $=$ Yosida &
$\nabla \delta_K^\varepsilon(u) = (N_K)_\varepsilon(u)
= \frac{1}{\varepsilon}(u - \psi)_-$ \\
Sistema penalizado &
$A u + \frac{1}{\varepsilon}(u-\psi)_- = b$ \\
Matriz de Newton &
$V_k = A + \frac{1}{\varepsilon}\operatorname{diag}(\chi_{\mathcal{A}_k})$,
\quad SPD \\
Convergencia & $u_\varepsilon \to u^\ast$, error $O(\sqrt{\varepsilon})$;
Newton cuadrático \\
\bottomrule
\end{tabular}
\end{center}

El hilo conductor: la indicatriz es la función convexa más brutal posible
(una pared infinita), su subdiferencial es el cono normal (un operador
multivaluado maximal monótono imposible de linearizar), y sin embargo la
regularización de Moreau--Yosida la convierte en un funcional diferenciable
cuyo gradiente es una función semisuave explícita, $(1/\varepsilon)
(u-\psi)_-$. Sobre esa función actúa Newton semisuave, cuya linealización
no es otra cosa que el método de conjunto activo. El análisis convexo
(subdiferencial, cono normal, resolvente) y el análisis numérico
(penalización, Newton, conjunto activo) son la misma historia contada
en dos idiomas.

% ============================================================
\begin{thebibliography}{9}

\bibitem{BS68}
H.~Brezis, G.~Stampacchia,
\emph{Sur la régularité de la solution d'inéquations elliptiques},
Bull. Soc. Math. France \textbf{96} (1968), 153--180.

\bibitem{Glowinski}
R.~Glowinski, J.-L.~Lions, R.~Trémolières,
\emph{Analyse numérique des inéquations variationnelles},
Dunod, Paris, 1976.

\bibitem{HIK02}
M.~Hinterm\"uller, K.~Ito, K.~Kunisch,
\emph{The primal-dual active set strategy as a semismooth Newton method},
SIAM J. Optim. \textbf{13} (2002), 865--888.

\bibitem{FacchineiPang}
F.~Facchinei, J.-S.~Pang,
\emph{Finite-Dimensional Variational Inequalities and Complementarity
Problems}, Springer, 2003.

\bibitem{Ulbrich}
M.~Ulbrich,
\emph{Semismooth Newton Methods for Variational Inequalities and
Constrained Optimization Problems in Function Spaces},
SIAM, 2011.

\bibitem{KinderlehrerStampacchia}
D.~Kinderlehrer, G.~Stampacchia,
\emph{An Introduction to Variational Inequalities and Their
Applications}, Academic Press, 1980.

\bibitem{Rockafellar70}
R.~T.~Rockafellar,
\emph{On the maximal monotonicity of subdifferential mappings},
Pacific J. Math. \textbf{33} (1970), 209--216.

\bibitem{Moreau65}
J.-J.~Moreau,
\emph{Proximité et dualité dans un espace hilbertien},
Bull. Soc. Math. France \textbf{93} (1965), 273--299.

\bibitem{Minty62}
G.~J.~Minty,
\emph{Monotone (nonlinear) operators in Hilbert space},
Duke Math. J. \textbf{29} (1962), 341--346.

\end{thebibliography}

\end{document}
